% R008 — Guide to Cultivating Complex Analysis (Bahasa Indonesia)
% Reader-facing translation unit: source/ca-v1.9/ca.tex lines 673–733.
% Source release: v1.9 (commit a4d25d9ba37a486a5a020219eb8bd4e6602fd14c).
% Translation provenance: OpenAI Codex gpt-5.6-sol, Ultra, acting on the user's request.
% Preserve source labels, glossary hooks, mathematical notation, and structure.

\chapter{Bidang Kompleks} \label{ch:complexplane}

\begin{myepigraph}
Jelas ini anggaran. Isinya banyak angka.

---George W.\ Bush
\end{myepigraph}

%%%%%%%%%%%%%%%%%%%%%%%%%%%%%%%%%%%%%%%%%%%%%%%%%%%%%%%%%%%%%%%%%%%%%%%%%%%%%%

\section{Bilangan kompleks} \label{sec:complexnums}

Matematika modern\footnote{Dalam konteks ini, \emph{modern} berarti
\myquote{setelah Abad Pertengahan.}}
bermula dari mengambil pernyataan yang salah seperti
\myquote{semua polinom memiliki akar} dan mendefinisikan ulang apa yang dapat
dimaksud dengan \myquote{akar}, yaitu mendefinisikan ulang \myquote{bilangan},
agar pernyataan tersebut menjadi benar.
Dalam hal ini, kita sampai pada bilangan kompleks.
Meskipun teknik ini (menggeser tiang gawang) terasa seperti curang, teknik ini
memberikan pada kita hampir seluruh matematika yang kita kenal, baik murni
maupun terapan.
Teknik yang sama dimulai dengan bilangan asli
\glsadd{not:N}$\N = \{ 1,2,3,\ldots \}$, satu-satunya
bilangan yang tampak langsung dari alam,
kemudian memberikan bilangan nol dan bilangan negatif sehingga menghasilkan
bilangan bulat \glsadd{not:Z}$\Z$, agar kita dapat menyelesaikan persamaan
seperti $n+2 = 1$.
Dari \glsadd{not:Z}$\Z$, kita mendefinisikan bilangan rasional
\glsadd{not:Q}$\Q$ untuk menyelesaikan\footnote{%
Apakah kita benar-benar menyelesaikan $2x=1$ dengan menulis
$x = \nicefrac{1}{2}$? Bagaimanapun, $\nicefrac{1}{2}$ hanyalah pengganti
untuk suatu objek yang tidak dapat kita uraikan dengan cara apa pun selain
\myquote{apa pun yang akan menjadi hasil $1$ dibagi $2$ seandainya objek itu
ada.}}
persamaan seperti $2x=1$.
Kita memperluas $\Q$ menjadi bilangan real \glsadd{not:R}$\R$
untuk menyelesaikan persamaan seperti $x^2=2$.
Sebenarnya, definisi bilangan real kita dibuat sedemikian rupa sehingga kita
mendapatkan teorema seperti teorema nilai antara, Bolzano--Weierstrass,
dan sebagainya.
Maka, tidaklah terlalu jauh untuk melakukan hal yang sama ketika berusaha
menyelesaikan $z^2+1=0$.
Sebagaimana pada bilangan real, konsekuensi menambahkan $\sqrt{-1}$ ke dalam
campuran ini jauh lebih mendalam daripada sekadar menemukan akar polinom.

Menariknya, sementara langkah menuju analisis dengan bilangan real merupakan
langkah menuju jurang, langkah menuju analisis dengan bilangan kompleks adalah
langkah menuju negeri ajaib dalam dongeng.
Mata kuliah analisis real tahun pertama menghancurkan harapan dan impian
mahasiswa.
Pernyataan yang masuk akal ternyata salah dan contoh tandingan yang aneh
berlimpah.
Sebaliknya, mata kuliah analisis kompleks mengisi mahasiswa dengan optimisme
yang tidak realistis.
Mata kuliah ini dipenuhi pernyataan naif dan konyol yang hanya mungkin
dianggap masuk akal oleh mahasiswa kalkulus yang buruk.\footnote{%
Misalnya: Jika Anda dapat mendiferensialkan sekali, Anda dapat
mendiferensialkan dua kali.
Setiap fungsi kurang lebih berperilaku seperti fungsi linear.
Jika semua turunan di suatu titik bernilai nol, fungsi tersebut konstan.
Dan seterusnya.}
Keduanya adalah polisi baik/polisi buruk, yin/yang,
dalam analisis kontemporer.
